% src/Prologo.tex — Prólogo (ensayo continuo) v0.2.x (2026-07-14)
% Reescritura completa a partir de prologo-geometria-de-la-informacion.md
% (entregado en la raíz del proyecto).
%
% Arco: pregunta (monedas 0.50/0.51 vs 0.98/0.99) → Ana/Beto
% (reparametrización) → truco (información de Fisher) → dónde está la
% geometría (variedad estadística) → por qué no es barniz (Chentsov + 5
% consecuencias verificables) → invitación al cierre.
%
% Caracteres transversales: Ana y Beto aparecen también en Cap.5, así
% que el lector los reencontrará como los analistas paradigmáticos de la
% invariancia.

\chapter*{Prólogo}
\addcontentsline{toc}{chapter}{Prólogo}

\begingroup
\setlength{\parindent}{0pt}
\setlength{\parskip}{7pt}
\emergencystretch=3em
\sloppy

\begin{flushright}
	\itshape \guillemotleft La incertidumbre tiene, además, una forma.\guillemotright
\end{flushright}

\vspace{1em}
Hay una pregunta que parece tonta hasta que se piensa dos veces: ¿qué tan lejos está una moneda que sale cara el 50\% de las veces de una que sale cara el 51\%? La respuesta obvia es ``un uno por ciento''. Pero esa respuesta es una trampa, y este libro existe para explicar por qué.

Imagina dos monedas. La primera tiene probabilidad 0.50 de cara, la segunda 0.51. Lánzalas cien veces cada una y, honestamente, casi no vas a poder distinguirlas: los resultados se van a mezclar, las dos van a dar rachas parecidas de caras y cruces. Ahora imagina otro par: una moneda con probabilidad 0.98 y otra con 0.99. También un uno por ciento de diferencia. Pero esta vez, con muy pocos lanzamientos, vas a notar la diferencia enseguida: la primera moneda te va a dar alguna cruz de vez en cuando, la segunda casi nunca. La misma resta ---$0.01$--- es un cambio minúsculo cerca del 50\% y un cambio enorme cerca del 100\%. Los números dicen que la diferencia es igual en ambos casos. Los datos dicen otra cosa.

Podrías pensar que esto es una curiosidad de juguete. No lo es. Es, literalmente, un problema que ya tuviste sobre tu escritorio la última vez que corriste una regresión logística.

\section*{Un problema que ya viviste sin saberlo}

Imagina dos analistas, Ana y Beto, ajustando la misma regresión logística con los mismos datos. Ana reporta todo en probabilidades: para su estimador $\widehat{\pi}$ da un error estándar de aproximadamente $\sqrt{\widehat{\pi}(1-\widehat{\pi})/n}$. Beto ajusta el mismo modelo pero reporta en la escala en la que el software realmente optimiza ---lo que se llama el \textbf{logit}, $\eta = \log\!\left(\dfrac{\pi}{1-\pi}\right)$---, y para $\widehat{\eta}$ obtiene un error estándar de aproximadamente $\dfrac{1}{\sqrt{\,n\,p(1-p)\,}}$.

Si nunca oíste la palabra ``logit'': es simplemente la transformación que usa \texttt{glm()} en R, \texttt{statsmodels} en Python, \texttt{PROC\ LOGISTIC} en SAS, y prácticamente todo software de regresión logística que hayas usado. La razón por la que existe es aburrida y necesaria: una probabilidad está encerrada entre 0 y 1, pero un coeficiente de regresión $\beta_0 + \beta_1 x$ puede dar cualquier número real. El logit ``desenrolla'' el intervalo $(0,1)$ a toda la recta para que la regresión lineal tenga sentido. Cada vez que reportaste un \emph{odds ratio} ---el pan de cada día en epidemiología, ciencias sociales, riesgo crediticio--- estabas exponenciando un coeficiente que vivía en logits.

Entonces: Ana y Beto ajustaron exactamente el mismo modelo, con los mismos datos, y llegan al mismo estimador puntual del odds ratio, a la misma verosimilitud, al mismo test de Wald hecho a su manera. Pero si les pides la cota de Cramér--Rao, o comparas directamente sus errores estándar, \textbf{los números no cuadran de forma trivial} --- uno es aproximadamente el recíproco del otro salvo constantes. Dos analistas honestos, mismo ajuste, reportando ``la incertidumbre del parámetro'' de maneras que no se traducen la una en la otra con una simple regla de tres.

Eso tiene nombre, y es el nombre que probablemente nunca cruzaste con estadística estimación: \textbf{reparametrización}. Reparametrizar es, sin más misterio, describir el mismo modelo usando una etiqueta distinta para el eje de parámetros --- $\pi$ en vez de $\eta$, o viceversa. Suena inofensivo, como cambiar de centímetros a pulgadas. Pero no lo es, porque centímetros y pulgadas están relacionados por una regla de conversión \emph{lineal} (multiplicar por una constante), y eso hace que las distancias se reescalen de forma consistente en todo el rango. La relación entre $\pi$ y $\eta = \log(\pi/(1-\pi))$ \textbf{no es lineal}: se estira muchísimo cerca de los extremos (0 y 1) y casi no hace nada cerca de 0.5. Es exactamente el mismo fenómeno que viste con las monedas de 0.98 y 0.99: la misma ``distancia en números'' significa cosas completamente distintas según dónde estés parado. Ana y Beto no cometieron ningún error. Cada uno usó una regla válida. El problema es que \textbf{ninguna de las dos reglas es la correcta}, porque ninguna es invariante ante cómo decidas nombrar al parámetro --- y ``la incertidumbre de mi modelo'' no debería depender de una decisión de notación.

\section*{El truco: dejar de medir con la regla equivocada}

Aquí es donde entra la idea central del libro.

Cuando estudiaste geometría, casi seguro empezaste con el plano de toda la vida: distancias con Pitágoras, líneas rectas como el camino más corto. Pero seguramente ya sabes, por geografía o por curiosidad, que esa no es la única geometría posible. Sobre la superficie de una esfera, el camino más corto entre Madrid y Buenos Aires no es una línea recta dibujada en un mapa plano --- es un arco curvado, porque la superficie misma está curvada. Un mapa plano de la Tierra (una proyección de Mercator, por ejemplo) distorsiona: cerca del ecuador es casi fiel, cerca de los polos exagera brutalmente las distancias, al punto de que Groenlandia parece del tamaño de África sin serlo ni de lejos. El mapa no cambia la Tierra; cambia cómo la medimos mal.

El espacio de los modelos estadísticos ---el espacio donde viven todos los valores posibles de $\pi$, o de $\eta$, o de cualquier parámetro--- tiene exactamente este problema. $\pi$ y $\eta$ son dos mapas distintos del mismo territorio, y ninguno de los dos es plano de verdad; solo Ana cree que el suyo lo es, y Beto cree que el suyo lo es, y ambos están viendo la misma distorsión desde ángulos distintos.

Lo que hace falta no es elegir un mapa ``mejor''. Hace falta dejar de medir distancias con el mapa y empezar a medirlas con el territorio. La herramienta que permite eso es una cantidad que primero hay que entender desde cero: \textbf{información}. En el lenguaje de todos los días ``información'' es casi sinónimo de ``datos''. Aquí va a significar algo más preciso: \textbf{información es cuánto te sorprende algo, medido en números.} Si un evento tiene probabilidad $p$, su información se define como $I = -\log(p)$: casi cero si $p$ está cerca de 1 (nada te sorprende), grande si $p$ está cerca de 0 (algo raro ocurrió). Esta cantidad, aplicada no a un evento sino a \emph{cuánto te ayuda un dato a distinguir un modelo de otro modelo vecino}, se llama \textbf{información de Fisher}. Y lo notable es que la información de Fisher, calculada correctamente, produce \textbf{la misma respuesta sin importar si trabajaste en $\pi$ o en $\eta$} --- porque no mide ``cuánto cambió el número del parámetro'', mide ``cuánto cambió realmente el modelo'', y esa segunda pregunta no depende de qué etiqueta le pusiste al eje. Convertida en una regla que cambia de tamaño punto a punto ---corta donde los datos son muy sensibles, larga donde no lo son--- la información de Fisher se convierte en lo que el cálculo multivariable y el álgebra lineal llaman una \textbf{métrica}: la única que hace que la distancia entre dos distribuciones no dependa de la parametrización que usaste para describirlas. Si conviertes correctamente el error estándar de Beto a la escala geométrica ---no con una resta ingenua, sino con la distancia inducida por Fisher--- coincide exactamente con el de Ana. La discrepancia desaparece. No por acuerdo entre ellos, sino porque había una sola respuesta correcta desde el principio.

\section*{¿Dónde está, literalmente, la geometría?}

Antes de seguir, hay que matar una confusión que si no se aclara ahora te va a perseguir todo el libro. Cuando escuchas ``geometría'' seguramente piensas en reglas, compases, curvas, la gráfica de una función $f(x)$ dibujada en un plano. Esa geometría ---la curvatura de una campana de Gauss dibujada en papel, medida con $f''(x)$--- \textbf{no tiene nada que ver con lo que sigue.} Es un objeto legítimo del cálculo de una variable, pero es ajeno por completo a la geometría de la información. Esa curvatura describe la forma de \emph{una} distribución. Aquí no nos interesa la forma de una distribución: nos interesa la relación entre \emph{varias}.

Imagina un archivero. Cada gaveta tiene una etiqueta $\theta$ ---digamos, un par $(\mu,\sigma)$--- y adentro de cada gaveta no hay un número: hay una campana de Gauss \emph{completa}, la función entera $f(x;\mu,\sigma)$ desplegada sobre todo el eje $x$. El archivero completo, el conjunto de todas las gavetas posibles indexadas por $\theta$, es el objeto que tiene geometría. Moverte de una gaveta a otra no es caminar sobre una curva; es cambiar de curva entera. Preguntar ``¿qué tan lejos está la gaveta $(\mu=0,\sigma=1)$ de la gaveta $(\mu=0.1,\sigma=1)$?'' no dice nada sobre valores de $x$: pregunta qué tan distinguibles son, como objetos completos, las dos distribuciones que hay adentro. Ese archivero es la variedad estadística. Vive en el espacio de parámetros $\Theta$ --- no en el espacio muestral donde vive $x$.

¿Y por qué hablamos desde el espacio de parámetros y no directamente sobre el conjunto de distribuciones que ese parámetro genera? Porque en realidad sí hablamos sobre ese conjunto --- pero ese conjunto es un subespacio de \emph{todas} las distribuciones posibles, un espacio de dimensión infinita, inmanejable de forma directa. Lo que hacemos es lo mismo que se hace con cualquier superficie curva en geometría diferencial clásica: en vez de manipular la esfera con sus coordenadas ambiente $(x,y,z)\in\mathbb{R}^3$, se usa una carta ---latitud y longitud--- y se ``jala hacia atrás'' la métrica ambiente a esa carta. $\Theta$ es esa carta. La familia de distribuciones $\{p_\theta\}$ es la superficie verdadera. Mientras la parametrización sea identificable ---dos valores distintos de $\theta$ dan distribuciones distintas---, trabajar en $\Theta$ con la métrica jalada hacia atrás es exactamente trabajar sobre la superficie, no una aproximación.

Y hay una razón técnica precisa por la que $x$ desaparece del todo. La información de Fisher se define como una \emph{esperanza} sobre $x$:
\[
	g_{ij}(\theta) \;=\; \mathbb{E}_{x\sim p_\theta}\!\left[\frac{\partial \log p_\theta(x)}{\partial \theta_i}\,\frac{\partial \log p_\theta(x)}{\partial \theta_j}\right].
\]
Esa $\mathbb{E}_x$ es una integral sobre todo el espacio muestral: $x$ se consume, se promedia, y lo que sobrevive depende únicamente de $\theta$. Eso no es casualidad, es diseño: la pregunta ``¿qué tan distinguibles son dos modelos?'' tiene que promediar sobre todos los datos que esos modelos podrían generar, no puede depender de un dato particular. Por eso la curvatura de este libro vive en el espacio de parámetros, y por eso Ana y Beto ---que discutían sobre $\pi$ y $\eta$, nunca sobre un valor concreto de $x$--- estaban, sin saberlo, discutiendo geografía del archivero.

\section*{Por qué esto no es formalismo vacío}

Este no es solo un vocabulario nuevo para renombrar cosas conocidas. Hay un resultado, el \textbf{teorema de Chentsov} (1972), que demuestra algo mucho más fuerte: si exiges que tu noción de distancia entre modelos sea invariante ante reparametrizaciones ---que no le importe si mides en $\pi$ o en $\eta$ o en cualquier otra escala---, entonces la métrica de Fisher es, salvo un factor de escala, \textbf{la única} que cumple esa exigencia. No es una entre varias opciones razonables. Es la única opción coherente. Del mismo tipo de resultado que dice que la distancia euclidiana es la única invariante ante rotaciones y traslaciones rígidas.

Y de ahí salen consecuencias que se pueden verificar, no solo admirar:

\textbf{Primero,} el caso de Ana y Beto deja de ser una anécdota incómoda y se convierte en una corrección calculable: sabes exactamente cómo traducir la incertidumbre reportada en una escala a la otra sin perder ni inventar información.

\textbf{Segundo,} cuando un algoritmo de aprendizaje ---una red neuronal, un modelo bayesiano--- avanza ``en la dirección de mayor mejora'' usando la resta ingenua de parámetros, está cometiendo el mismo error que Ana o Beto por separado: moverse según el mapa, no según el territorio. Corregir el paso con la métrica de Fisher (el llamado \emph{gradiente natural}) hace que el algoritmo avance según cuánto cambia realmente el modelo, y eso cambia de forma medible la velocidad y estabilidad de la convergencia.

\textbf{Tercero,} hay familias de distribuciones donde este espacio resulta ser, en un sentido preciso, plano ---y ahí la intuición ingenua de Ana o de Beto casi funciona sin corrección---. Hay otras donde el espacio está fuertemente curvado, y ahí la intuición va a fallar sistemáticamente. Saber \emph{dónde} y \emph{cuánto} va a fallar es, en sí mismo, información valiosa sobre los límites de la inferencia estadística clásica.

\textbf{Cuarto, el mismo problema de Ana y Beto contamina tus creencias previas antes de ver un solo dato.} En estadística bayesiana es tentador poner un prior ``sin preferencias'', uniforme en $\pi\in(0,1)$. Pero un prior uniforme en $\pi$ no es el mismo prior que uno uniforme en $\eta$: al cambiar de coordenada la densidad se distorsiona por un jacobiano, así que ``no saber nada'' en una escala resulta ser una creencia muy fuerte en la otra. Jeffreys (1946) resolvió esto definiendo el prior como $\pi(\theta)\propto\sqrt{\det g(\theta)}$ ---el elemento de volumen de la misma variedad riemanniana que ya construiste---, que es, otra vez por el teorema de Chentsov, la única elección que no depende de en qué coordenadas decidiste ser ignorante.

Y quinto ---quizás el más convincente si ya tienes años haciendo estadística aplicada---, \textbf{esto no es nuevo: ya lo usaste sin el nombre.} La transformación $\arcsin(\sqrt{\hat p})$ para estabilizar varianza de proporciones, el $\sqrt{x}$ para conteos de Poisson, la $z$ de Fisher para correlaciones --- todas esas recetas clásicas para ``que el modelo lineal aplique'' se derivan pidiendo que $h'(\theta)\propto\sqrt{I(\theta)}$, es decir, calculando $h(\theta)=\int\sqrt{I(\theta)}\,d\theta$. Eso es exactamente la longitud de arco sobre la métrica de Fisher en una variedad de una sola dimensión: la coordenada donde el archivero por fin tiene gavetas parejas y una regla ordinaria vuelve a servir. Con un solo parámetro esa coordenada mágica siempre existe, por eso varias generaciones de estadísticos resolvieron el problema con una tabla de transformaciones sin necesitar el aparato completo de Riemann. En cuanto el modelo tiene dos o más parámetros, esa escapatoria generalmente desaparece ---la curvatura deja de ser eliminable con un solo cambio de variable--- y ahí es donde el truco de una dimensión se queda corto y hace falta la teoría completa.

\section*{La invitación}

No te voy a prometer que este libro es fácil, pero sí que no vas a leer una sola idea que no puedas anclar en algo que ya conocías: un glm que ya corriste, un odds ratio que ya reportaste, una discrepancia entre errores estándar que quizás ya notaste y le echaste la culpa a la aproximación asintótica. De ahí en más vamos a apoyarnos en tu álgebra lineal, tu cálculo multivariable y tu estadística matemática para darle forma rigurosa a esa intuición. Lo que te prometo es que, cuando termines, nunca más vas a mirar un espacio de parámetros como una simple lista de números en una escala arbitraria. Vas a ver una superficie con forma propia, con caminos más cortos que no son líneas rectas, con regiones donde dos modelos casi no se distinguen entre sí aunque sus números digan lo contrario.

La estadística te enseñó a hacer preguntas sobre la incertidumbre. La geometría de la información te va a enseñar que esa incertidumbre tiene, además, una \emph{forma} --- y que esa forma es lo único en lo que Ana y Beto, sin saberlo, ya estaban de acuerdo.

\vspace{2em}

\begin{flushright}
	\itshape \guillemotleft Empecemos por medir bien.\guillemotright
\end{flushright}

\endgroup
